\chapter{Liver}

\label{ch:ch45_liver}

The liver is the largest internal organ and performs over 500 functions. Imaging of the liver helps detect tumors, cysts, fatty changes, scarring (cirrhosis), infections, and blood flow problems. The three main imaging tools for the liver are:

\begin{itemize}
  \item \textbf{Ultrasound}: Often the first test; quick, safe, no radiation. Good for detecting fatty liver, large masses, bile duct dilation, and assessing blood flow with Doppler. Limited by body habitus and operator skill.
  \item \textbf{CT with contrast}: Excellent for detecting and characterising liver lesions. A multiphasic protocol (arterial, portal venous, and delayed phases) is essential because different liver lesions show characteristic enhancement patterns at different time points.
  \item \textbf{MRI}: Provides the most detailed characterisation of liver lesions. Certain sequences (T2-weighted, diffusion-weighted, hepatobiliary phase after liver-specific contrast) can diagnose benign lesions such as haemangiomas and focal nodular hyperplasia (FNH) with near-100\% confidence, often avoiding the need for biopsy.
\end{itemize}

Your doctor will choose the right test based on the clinical question, your liver function, and whether you are at risk for liver cancer (e.g., due to cirrhosis or hepatitis).


\section{Imaging}
Ultrasound is often a first study for liver and biliary symptoms; multiphasic CT or MRI characterizes focal lesions and vascularity. Contrast-enhanced ultrasound or nuclear imaging may answer selected questions.

\begin{table}[h]
\centering
\caption{Multiphasic CT/MRI protocol for the liver}
\begin{tabular}{p{3cm} p{3.5cm} p{5cm}}
\hline
\textbf{Phase} & \textbf{Timing post-injection} & \textbf{What it evaluates} \\
\hline
Arterial (late arterial / hepatic arterial) & 20--35 s & HCC arterial enhancement, FNH (intense enhancement), hypervascular metastases \\
Equilibrium / delayed & 3--5 min & Hemangioma (peripheral nodular enhancement with centripetal fill-in); washout of HCC; capsular enhancement \\
Hepatobiliary phase (MRI only) & 15--20 min (gadoxetate) & FNH retention, well-differentiated HCC retention, metastases (hypointense) \\
Non-contrast & Pre-injection & Calcification, fat, acute hemorrhage \\
Portal venous & 60--70 s & Standard phase; most liver lesions appear hypoattenuating \\
\hline
\end{tabular}
\end{table}

\begin{pearlbox}[Multiphase protocol is essential]
A single-phase CT or MRI is insufficient to characterize liver lesions. The arterial phase detects hypervascular lesions; the portal venous phase detects hypovascular lesions and provides the best parenchymal assessment; the delayed phase evaluates washout and enhancement patterns. Always request multiphase imaging for liver lesion characterisation.
\end{pearlbox}


\section{Focal Liver Lesions}

\subsection{Benign Lesions}

\begin{table}[h]
\centering
\caption{Benign liver lesions: imaging features}
\begin{tabular}{p{2.5cm} p{3.5cm} p{3.5cm} p{3cm}}
\hline
\textbf{Lesion} & \textbf{CT} & \textbf{MRI} & \textbf{Key feature} \\
\hline
FNH & Intense arterial enhancement; iso/hypoattenuating on portal venous; central scar (hyperdense on delayed) & Isointense to liver on T1/T2; enhancement in hepatobiliary phase (contains hepatocytes); central scar T2 bright & Usually in young women; no treatment needed \\
Hemangioma & Peripheral nodular enhancement with centripetal fill-in; remains hyperattenuating on delayed phases & T2 very bright (``light bulb sign''); same enhancement pattern as CT & Most common benign liver tumor \\
Hepatic adenoma & Homogeneous arterial enhancement; may contain fat & T1 hyperintense (fat); no hepatobiliary phase uptake & Risk of hemorrhage and malignant transformation; associated with OCP use \\
Simple cyst & Water density ($-10$ HU); no enhancement & T1 very dark, T2 very bright; no enhancement & No treatment needed \\
\hline
\end{tabular}
\end{table}

\subsection{Malignant Lesions}

\begin{table}[h]
\centering
\caption{Malignant liver lesions: imaging features}
\begin{tabular}{p{2.5cm} p{3.5cm} p{3.5cm} p{3cm}}
\hline
\textbf{Lesion} & \textbf{CT} & \textbf{MRI} & \textbf{Key feature} \\
\hline
Cholangiocarcinoma & Peripheral enhancement with progressive centripetal fill-in; capsular retraction & T2 hyperintense; delayed enhancement; capsular retraction & Arises from bile duct epithelium \\
HCC & Arterial enhancement with portal venous/delayed washout; capsule; threshold growth & T1 variable; T2 mildly hyperintense; washout on dynamic; hepatobiliary phase hypointense & Most common primary liver cancer; arises in cirrhosis \\
Metastases & Typically hypovascular (hypoattenuating on portal venous); rim enhancement for some & T2 mildly hyperintense; diffusion restriction; hypointense on hepatobiliary phase & Most common liver malignancy overall; often multiple \\
\hline
\end{tabular}
\end{table}

\begin{pearlbox}[Hemangioma vs metastasis]
Both can appear as round enhancing liver lesions. Haemangiomas show \textbf{peripheral nodular enhancement that fills in centrally over time} and are very bright on T2. Metastases typically show \textbf{rim enhancement} and do NOT fill in. If uncertain, MRI with hepatobiliary phase (gadoxetate) can confirm hemangioma with near-100\% confidence.
\end{pearlbox}


\section{LI-RADS}
LI-RADS (Liver Imaging Reporting and Data System) standardises the reporting of liver lesions in patients at risk for HCC (cirrhosis, chronic hepatitis B, prior HCC). It applies ONLY to multiphasic CT or MRI with appropriate protocol.

\begin{table}[h]
\centering
\caption{LI-RADS categories}
\begin{tabular}{p{2cm} p{3.5cm} p{3cm} p{3cm}}
\hline
\textbf{Category} & \textbf{Description} & \textbf{Major features} & \textbf{Management} \\
\hline
LR-1 & Definitely benign & Cyst, hemangioma, etc. & No further action \\
LR-2 & Probably benign & Small, non-specific & Short-term follow-up \\
LR-3 & Intermediate probability & Some HCC features but not all & Follow-up \\
LR-4 & Probably HCC & Arterial hyperenhancement + one major feature & Treat or biopsy \\
LR-5 & Definitely HCC & Arterial hyperenhancement + $\geq 2$ major features & Treat without biopsy \\
LR-M & Probably malignant, not HCC & Non-HCC features (e.g., mass-forming, rim enhancement) & Biopsy \\
LR-TIV & Tumor in vein & Enhancing tumor in portal/hepatic vein & Treat \\
\hline
\end{tabular}
\end{table}

\textbf{LI-RADS major features for HCC}:
\begin{enumerate}
  \item Arterial phase hyperenhancement (APHE)
  \item Size $\geq 10$ mm
  \item Washout (portal venous or delayed phase)
  \item Enhancing capsule (delayed phase)
  \item Threshold growth ($\geq 50\%$ size increase in $\leq 6$ months)
\end{enumerate}

\begin{pearlbox}[LI-RADS 5 = HCC, no biopsy needed]
In a patient with cirrhosis, an LI-RADS 5 lesion (arterial hyperenhancement + washout + capsule + size $\geq 10$ mm, or threshold growth) can be diagnosed as HCC and treated without histological confirmation. This is an exception to the general rule that tumors require biopsy for diagnosis.
\end{pearlbox}


\section{Diffuse Liver Disease}

\subsection{Fatty Liver (Hepatic Steatosis)}
\begin{itemize}
  \item \textbf{CT}: Liver parenchyma is less dense than the spleen (liver attenuation $<$ splenic attenuation on unenhanced CT; $< 40$ HU suggests significant steatosis).
  \item \textbf{MRI}: Proton density fat fraction (PDFF) quantifies steatosis non-invasively; $> 5\%$ is abnormal.
  \item \textbf{Focal sparing}: Typically in the posterior segments of the right lobe and around the falciform ligament; can mimic a mass.
\end{itemize}

\subsection{Cirrhosis}
Imaging features of cirrhosis:
\begin{itemize}
  \item \textbf{Morphology}: Nodular liver surface, caudate lobe hypertrophy, atrophy of right and left lobes
  \item \textbf{Parenchyma}: Heterogeneous enhancement, regenerative and dysplastic nodules
  \item \textbf{Portal hypertension}: Splenomegaly ($> 13$ cm), portosystemic varices, ascites, recanalised umbilical vein
  \item \textbf{Complications}: HCC surveillance required (LI-RADS), variceal bleeding, hepatic encephalopathy
\end{itemize}


\section{Portal Hypertension Imaging}
\begin{table}[h]
\centering
\caption{Signs of portal hypertension on imaging}
\begin{tabular}{p{5cm} p{5cm}}
\hline
\textbf{Finding} & \textbf{Significance} \\
\hline
Ascites & May indicate decompensated cirrhosis \\
Hepatofugal flow (reversed portal flow) & Severe portal hypertension \\
Portal vein diameter $> 13$ mm & Suggests portal hypertension \\
Portosystemic varices & Oesophageal, gastric, splenorenal, umbilical \\
Recanalised paraumbilical vein & Persistent portal hypertension \\
Splenomegaly $> 13$ cm & Chronic portal hypertension \\
\hline
\end{tabular}
\end{table}


\section{Liver Transplantation Imaging}
Pre-transplant evaluation includes:
\begin{itemize}
  \item \textbf{Vascular anatomy}: Hepatic artery variant anatomy (replaced right or left hepatic artery in 20--40\% of the population) is critical for surgical planning---document on CTA or MRA.
  \item \textbf{Tumor burden}: Milan criteria (single tumor $\leq 5$ cm or $\leq 3$ tumors each $\leq 3$ cm, no vascular invasion, no extrahepatic disease) determines transplant eligibility.
  \item \textbf{Portal vein patency}: Thrombosis may alter surgical approach.
\end{itemize}

Post-transplant complications assessed by imaging:
\begin{itemize}
  \item \textbf{Hepatic artery thrombosis}: Most common vascular complication; graft ischemia
  \item \textbf{Portal vein stenosis}: Narrowing at the anastomosis
  \item \textbf{Biliary complications}: Stricture, leak (``ischaic cholangiopathy'' in arterial thrombosis)
\end{itemize}


\section{Interpretation}
Assess liver contour, background parenchyma, portal and hepatic veins, biliary dilation, focal lesion size and enhancement, fat, iron, and signs of portal hypertension. Use LI-RADS only in the population and protocol for which it is intended.

\section{Safety}
Contrast selection, kidney function, pregnancy, and prior reaction require individualized review. Report suspected abscess, acute vascular complication, biliary obstruction, or tumor thrombus promptly.

\clearpage
